\chapter{Of Amendment and the Continuance of the Pact}
\label{art:amendment}

\articlenote{In which is set down that this Pact, though the law that governs the silo, may be altered by the silo's citizens acting together; that no amendment may diminish the rights and protections this Pact gives to every citizen, nor contradict the fundamental necessity of the silo's survival and the equality of its citizens before the law; and that certain provisions stand beyond the Assembly's reach altogether, changeable by no vote and no petition, that the silo's continuance not be undone by well-meant hands.}

\section{The Process of Amendment}

This Pact may be amended by the Assembly, gathered in Assembly under Article 20, upon proposal from the Mayor or petition from one fifth of the silo's citizens of majority. An amendment requires the affirmative vote of two thirds of the Assembly present.

\begin{clauses}
\item The text of a proposed amendment shall be made known to all citizens not fewer than fourteen days before the Assembly vote, and citizens shall have the right to speak to the amendment in the Assembly or to send messages under Article 21 raising objections or support.
\item An amendment, once ratified, is entered in the Amendment Log maintained under Article 21 and becomes part of this Pact from the day of ratification forward. All earlier editions of this Pact remain in the archive for the record.
\item No amendment may be repealed save by a later amendment, ratified in the same manner as the first.
\end{clauses}

\section{The Limits on Amendment}

Certain truths are so fundamental to the silo's nature that no amendment may alter them, lest the silo cease to be a place where citizens may live.

\begin{clauses}
\item No amendment may eliminate a citizen's right to petition for a final request to go outside, nor may an amendment make that request answerable as an offense under this Pact.
\item No amendment may eliminate the lottery for children under Article 3, nor may an amendment increase a citizen's burden of labor beyond the capacity any ordinary citizen might bear.
\item No amendment may remove Judicial's right to rule where this Pact is silent, nor may an amendment subordinate Judicial to the Mayor's authority without the Assembly's explicit consent.
\item No amendment may eliminate a citizen's right to know the content of this Pact, to view records concerning the citizen under Article 21, or to speak before the Assembly under Article 20.
\item No amendment may codify inequality before the law; every citizen shall remain equal under the Pact, save where the Pact itself explicitly provides differences by office, trade, or age.
\item No amendment may alter the founding as Article 1 records it, nor compel any Department to disclose the manner, location, or working of any matter this Pact elsewhere commits to that Department's own procedure and does not itself set down.
\item An amendment touching any limit named in this section takes no effect upon the Assembly's vote alone. It further requires the written concurrence of the Judge and the Head of Information Technology, given only upon their joint finding that the amendment works none of the harms this section forbids. An amendment lacking such concurrence is void, whatever majority approved it, and is struck from the Log as though never ratified.
\end{clauses}

\section{The Amendment Log}

Amendments to this Pact, ratified by the Assembly, are recorded in a separate roll under Article 21 and are printed in the backmatter of this edition of the Pact, under Appendix A, so that any citizen may see how the law has grown and changed.

\begin{clauses}
\item Each amendment is entered with the date of ratification, the text of the amendment in full, and a brief notation of what Article or Section it amends or adds.
\item The Amendment Log is maintained in the Judicial archive, and a backup copy is kept by the Department of Supply under Article 11.
\item Any citizen may petition Judicial to examine the Amendment Log and to ask for an explanation of any amendment's import or consequence.
\end{clauses}

\section{The Continuance and the Renewal of the Pact}

This Pact is made for the continuance of the silo and the protection of every citizen who dwells here. It shall endure as long as the silo stands, and shall bind every citizen born or entering the silo until the citizen's death or final release.

\begin{clauses}
\item Every citizen of majority takes an oath to know this Pact, to obey it, and to call others to obey it, under Article 2. This oath is binding and sacred.
\item The Mayor, the Judge, the Sheriff, and every Head of Department swear at their assumption of office to defend this Pact against those who would breach it, and to enforce this Pact fairly and without favor.
\item Should the silo ever cease to stand, or should the citizens ever no longer dwell here, this Pact ends. But while the silo stands and citizens dwell within it, this Pact shall govern, and every citizen shall be bound by it.
\end{clauses}
